\documentclass[11pt, a4paper]{article}
\usepackage[utf8]{inputenc}
\usepackage[T1]{fontenc}
\usepackage[ngerman]{babel}
\usepackage{geometry}
\usepackage{helvet}
\usepackage{tabularx}
\usepackage{parskip}

% Seitenränder anpassen
\geometry{left=1.5cm, right=1.5cm, top=2cm, bottom=2cm}

% Schriftartwechsel auf Sans-Serif
\renewcommand{\familydefault}{\sfdefault}

\begin{document}

\section*{Project Canvas für die Masterarbeit}
\textbf{Datum:} \today \quad | \quad \textbf{Bearbeiter:} Moritz Diez

\vspace{0.5cm}

\begin{tabularx}{\textwidth}{|X|X|X|}
\hline
\textbf{ZWECK} & \textbf{BUDGET} & \textbf{TEAM} \\ 
\footnotesize
Ein KI-gestützter Ansatz zur Lösung des 'Schuhschachtel-Dilemmas' wird entwickelt. Unstrukturierte Bestandsdaten werden in modulare SysML v2 Modelle überführt. Implizites Wissen wird durch LLMs antizipiert, um die Nutzbarkeit der Modelle sicherzustellen. & 
\footnotesize
Kein explizites Budget definiert. (Standardressourcen der Hochschule/Firma). & 
\footnotesize
\textbullet\ Moritz Diez (Masterand) \newline
\textbullet\ [Betreuer HM] \newline
\textbullet\ [Betreuer Extern] \\ \hline

\textbf{UMFELD} & \textbf{MEILENSTEINE} & \textbf{QUALITÄT} \\ 
\footnotesize
Brownfield-Engineering, Transformation zu MBSE. Nutzung des technologischen Enablers SysML v2 für 'Everything as Code'. & 
\footnotesize
1. Konzept \& Analyse (15.04.) \newline
2. Prototyp Implementierung (31.05.) \newline
3. Validierung \& Test (15.07.) \newline
4. Abgabe (15.09.) & 
\footnotesize
Syntaktische Korrektheit des SysML v2 Codes. Gewährleistung der 'Plug \& Play'-Fähigkeit generierter Teilmodelle ohne Refaktorierung. \\ \hline

\textbf{RISIKEN} & \textbf{ERGEBNIS} & \textbf{KUNDE} \\ 
\footnotesize
\textbullet\ Fehlinterpretation der LLMs bei komplexen technischen Zusammenhängen. \newline
\textbullet\ Mangelnde Kompatibilität isoliert generierter Teilmodelle. & 
\footnotesize
Ein dreistufiger 'AiC-Generator' (Ingestion, Pattern Matching, Synthesis). Nachweis der Eignung von SysML v2 als KI-Enabler. & 
\footnotesize
Unternehmen mit Bestandsanlagen. Systems Engineering Community. \\ \hline
\end{tabularx}

\vspace{0.5cm}
\textbf{ZEITRAHMEN:} 15.03.2026 -- 15.09.2026

\end{document}